% vim: set textwidth=120:

% Example CV based on the 1.5-column-cv template. Main features:
% * uses the Roboto font family and IcoMoon icon set;
% * doesn't use colours, different font weights are used instead for styling;
% * because the CV fits on one page, header and footer is empty, since there isn't much useful info to put there;
% * includes a photo.
\documentclass[a4paper,10pt]{article}




% package imports
% ---------------

\usepackage[british]{babel} % for correct language and hyphenation and stuff
\usepackage{calc}           % for easier length calculations (infix notation)
\usepackage{enumitem}       % for configuring list environments
\usepackage{fancyhdr}       % for setting header and footer
\usepackage{fontspec}       % for fonts
\usepackage{geometry}       % for setting margins (\newgeometry)
\usepackage{graphicx}       % for pictures
\usepackage{microtype}      % for microtypography stuff
\usepackage{xcolor}         % for colours
\usepackage{hyperref}       % for links


% margin and column widths
% ------------------------

% margins
\newgeometry{left=15mm,right=15mm,top=15mm,bottom=15mm}

% width of the gap between left and right column
\newlength{\cvcolumngapwidth}
\setlength{\cvcolumngapwidth}{3.5mm}

% left column width
\newlength{\cvleftcolumnwidth}
\setlength{\cvleftcolumnwidth}{44.5mm}

% right column width
\newlength{\cvrightcolumnwidth}
\setlength{\cvrightcolumnwidth}{\textwidth-\cvleftcolumnwidth-\cvcolumngapwidth}

% set paragraph indentation to 0, because it screws up the whole layout otherwise
\setlength{\parindent}{0mm}


% style definitions
% -----------------
% style categories explanation:
% * \cvnameXXX is used for the name;
% * \cvsectionXXX is used for section names (left column, accompanied by a horizontal rule);
% * \cvtitleXXX is used for job/education titles (right column);
% * \cvdurationXXX is used for job/education durations (left column);
% * \cvheadingXXX is used for headings (left column);
% * \cvmainXXX (and \setmainfont) is used for main text;
% * \cvruleXXX is used for the horizontal rules denoting sections.

% font families
\defaultfontfeatures{Ligatures=TeX} % reportedly a good idea, see https://tex.stackexchange.com/a/37251

\newfontfamily{\cvnamefont}{Roboto Medium}
\newfontfamily{\cvsectionfont}{Roboto Medium}
\newfontfamily{\cvtitlefont}{Roboto Regular}
\newfontfamily{\cvdurationfont}{Roboto Light Italic}
\newfontfamily{\cvheadingfont}{Roboto Regular}
\setmainfont{Roboto Light}

% colours
\definecolor{cvnamecolor}{HTML}{000000}
\definecolor{cvsectioncolor}{HTML}{000000}
\definecolor{cvtitlecolor}{HTML}{000000}
\definecolor{cvdurationcolor}{HTML}{000000}
\definecolor{cvheadingcolor}{HTML}{000000}
\definecolor{cvmaincolor}{HTML}{000000}
\definecolor{cvrulecolor}{HTML}{000000}

\color{cvmaincolor}

% styles
\newcommand{\cvnamestyle}[1]{{\Large\cvnamefont\textcolor{cvnamecolor}{#1}}}
\newcommand{\cvsectionstyle}[1]{{\normalsize\cvsectionfont\textcolor{cvsectioncolor}{#1}}}
\newcommand{\cvtitlestyle}[1]{{\large\cvtitlefont\textcolor{cvtitlecolor}{#1}}}
\newcommand{\cvdurationstyle}[1]{{\small\cvdurationfont\textcolor{cvdurationcolor}{#1}}}
\newcommand{\cvheadingstyle}[1]{{\normalsize\cvheadingfont\textcolor{cvheadingcolor}{#1}}}


% inter-item spacing
% ------------------

% vertical space after personal info and standard CV items
\newlength{\cvafteritemskipamount}
\setlength{\cvafteritemskipamount}{5mm plus 1.25mm minus 1.25mm}

% vertical space after sections
\newlength{\cvaftersectionskipamount}
\setlength{\cvaftersectionskipamount}{2mm plus 0.5mm minus 0.5mm}

% extra vertical space to be used when a section starts with an item with a heading (e.g. in the skills section),
% so that the heading does not follow the section name too closely
\newlength{\cvbetweensectionandheadingextraskipamount}
\setlength{\cvbetweensectionandheadingextraskipamount}{1mm plus 0.25mm minus 0.25mm}


% intra-item spacing
% ------------------

% vertical space after name
\newlength{\cvafternameskipamount}
\setlength{\cvafternameskipamount}{3mm plus 0.75mm minus 0.75mm}

% vertical space after personal info lines
\newlength{\cvafterpersonalinfolineskipamount}
\setlength{\cvafterpersonalinfolineskipamount}{2mm plus 0.5mm minus 0.5mm}

% vertical space after titles
\newlength{\cvaftertitleskipamount}
\setlength{\cvaftertitleskipamount}{1mm plus 0.25mm minus 0.25mm}

% value to be used as parskip in right column of CV items and itemsep in lists (same for both, for consistency)
\newlength{\cvparskip}
\setlength{\cvparskip}{0.5mm plus 0.125mm minus 0.125mm}

% set global list configuration (use parskip as itemsep, and no separation otherwise)
\setlist{parsep=0mm,topsep=0mm,partopsep=0mm,itemsep=\cvparskip}


% CV commands
% -----------

% creates a "personal info" CV item with the given left and right column contents, with appropriate vertical space after
% @param #1 left column content (should be the CV photo)
% @param #2 right column content (should be the name and personal info)
\newcommand{\cvpersonalinfo}[3]{
    % left and right column
    \begin{minipage}[t]{\cvleftcolumnwidth}
        \vspace{0mm} % XXX hack to align to top, see https://tex.stackexchange.com/a/11632
        \raggedleft #1
    \end{minipage}% XXX necessary comment to avoid unwanted space
    \hspace{\cvcolumngapwidth}% XXX necessary comment to avoid unwanted space
    % \begin{minipage}[t]{\cvrightcolumnwidth}
    \begin{minipage}[t]{60mm}
        \vspace{0mm} % XXX hack to align to top, see https://tex.stackexchange.com/a/11632
        #2
    \end{minipage}
    % \begin{minipage}[t]{\cvrightcolumnwidth}
    \begin{minipage}[t]{70mm}
        \vspace{0mm} % XXX hack to align to top, see https://tex.stackexchange.com/a/11632
	#3
    \end{minipage}

    % space after
    \vspace{\cvafteritemskipamount}
}

% typesets a name, with appropriate vertical space after
% @param #1 name text
\newcommand{\cvname}[1]{
    % name
    \cvnamestyle{#1}

    % space after
    \vspace{\cvafternameskipamount}
}

% typesets a line of personal info beginning with an icon, with appropriate vertical space after
% @param #1 parameters for the \includegraphics command used to include the icon
% @param #2 icon filename
% @param #3 line text
\newcommand{\cvpersonalinfolinewithicon}[3]{
    % icon, vertically aligned with text (see https://tex.stackexchange.com/a/129463)
    \raisebox{.5\fontcharht\font`E-.5\height}{\includegraphics[#1]{#2}}
    % text
    #3

    % space after
    \vspace{\cvafterpersonalinfolineskipamount}
}

% creates a "section" CV item with the given left column content, a horizontal rule in the right column, and with
% appropriate vertical space after
% @param #1 left column content (should be the section name)
\newcommand{\cvsection}[1]{
    % left and right column
    \begin{minipage}[t]{\cvleftcolumnwidth}
        \raggedleft\cvsectionstyle{#1}
    \end{minipage}% XXX necessary comment to avoid unwanted space
    \hspace{\cvcolumngapwidth}% XXX necessary comment to avoid unwanted space
    \begin{minipage}[t]{\cvrightcolumnwidth}
        \textcolor{cvrulecolor}{\rule{\cvrightcolumnwidth}{0.3mm}}
    \end{minipage}

    % space after
    \vspace{\cvaftersectionskipamount}
}

% creates a standard, multi-purpose CV item with the given left and right column contents, parskip set to cvparskip
% in the right column, and with appropriate vertical space after
% @param #1 left column content
% @param #2 right column content
\newcommand{\cvitem}[2]{
    % left and right column
    \begin{minipage}[t]{\cvleftcolumnwidth}
        \raggedleft #1
    \end{minipage}% XXX necessary comment to avoid unwanted space
    \hspace{\cvcolumngapwidth}% XXX necessary comment to avoid unwanted space
    \begin{minipage}[t]{\cvrightcolumnwidth}
        \setlength{\parskip}{\cvparskip} #2
    \end{minipage}

    % space after
    \vspace{\cvafteritemskipamount}
}

% typesets a title, with appropriate vertical space after
% @param #1 title text
\newcommand{\cvtitle}[1]{
    % title
    \cvtitlestyle{#1}

    % space after
    \vspace{\cvaftertitleskipamount}
    % XXX need to subtract cvparskip here, because it is automatically inserted after the title "paragraph"
    \vspace{-\cvparskip}
}


% header and footer
% -----------------

% set empty header and footer
\pagestyle{empty}



% preamble end/document start
% ===========================

\begin{document}


% personal info
% -------------

\cvpersonalinfo{
    % photo
    \includegraphics[height=24mm]{me-circle.png}

}{
    % name
    \cvname{Timofei Tipishev}

    % email address
    \cvpersonalinfolinewithicon{height=4mm}{070-envelop.pdf}{
        tipishev@gmail.com
    }


    % phone number
    \cvpersonalinfolinewithicon{height=4mm}{067-phone.pdf}{
        +46 762\,333\,946
    }

    % GitHub account
    \cvpersonalinfolinewithicon{height=4mm}{github.png}{
	    \href{https://github.com/tipishev}{tipishev}
    }
}{
	\begin{itemize}[leftmargin=*]
		\item Data Scientist/Engineer/Plumber
		\item Senior Python/Django/Flask developer
		\item Tech workshops speaker
		\item JavaScript/Vue.js amateur
		\item Monkey C / Erlang / Haskell hobbyist
	\end{itemize}
}


% work experience
% ---------------

\cvsection{WORK EXPERIENCE}

% Sportamore
\cvitem{
    \cvdurationstyle{2016 -- present}
}{
    \cvtitle{Senior Software Developer}

    Sportamore

    % greenfield
    \begin{itemize}[leftmargin=*]
        \item Created and launched internal services
	    \begin{itemize}[leftmargin=*]
		  \item supplier orders recommendation service (Flask, scikit-learn, Vue.js)
		  \item reporting workflows runner (Apache Airflow, Google Sheets)
		  \item drag-and-drop promotional emails creator (Django, Vanilla.js)
		  \item product prices control center (Django Rest Framework)
	    \end{itemize}
    \end{itemize}

    % brownfield
    \begin{itemize}[leftmargin=*]
        \item Maintained and improved internal services 
	    \begin{itemize}[leftmargin=*]
		  \item data warehouse (Python 3, Flask, SQLAlchemy, RabbitMQ)
		  \item physical warehouse management system (Django 1.8)
		  \item the main website (Python 2, Django 1.11)
		  \item deployment server (Fabric, custom blueprints)
	    \end{itemize}
    \end{itemize}

    % integrations
    \begin{itemize}[leftmargin=*]
        \item Integrated with external systems
	    \begin{itemize}[leftmargin=*]
		  \item Klarna: handling payments and returns
		  \item Swisslog: managing robot-warehouse
		  \item SendGrid: scheduling emails
		  \item Unifaun: dispatching deliveries to customers
	    \end{itemize}
    \end{itemize}

    % workshops

    \begin{itemize}[leftmargin=*]
        \item Performed internal and external workshops
	    \begin{itemize}[leftmargin=*]
	          \item \textit{Replacing cron jobs with Airflow}
		  \item \textit{\href{https://github.com/tipishev/python_debugging_workshop}{Python debugging as a text adventure} (PyCon SE 2019)}
		  \item \textit{\href{https://github.com/tipishev/git_workshop}{Git tips and tricks}}
		  \item service launch presentations
	    \end{itemize}
    \end{itemize}

}

% Paxport
\cvitem{
    \cvdurationstyle{2015 -- 2016}
}{
    \cvtitle{Java Backend Developer}

    Paxport (as a Fastdev consultant)

    \begin{itemize}[leftmargin=*]
	\item Developed airplane seats booking backend using Java8/PostgreSQL/Spring/Vert.x
	\item Integrated booking backend with J2EE-monolith using Vert.x (eventBus/HTTP)
	\item Developed and supported Java-applications within the monolith
    \end{itemize}
}

% Yandex
\cvitem{
    \cvdurationstyle{2014 -- 2015}
}{
    \cvtitle{Python/Flask Backend Developer}

    Yandex

    \begin{itemize}[leftmargin=*]
	\item Developed and launched a Python/Flask/MongoDB/Celery web service
	\item Extended API functionality
	\item Integrated with internal services Yandex.Music, Yandex.News, etc.
	\item Created data importers from XML, JSON, and custom APIs
    \end{itemize}
}


% Education
% ---------

\cvsection{EDUCATION}

% coursera
\cvitem{
    \cvdurationstyle{2020 -- present}
}{
    \cvtitle{Reinforcement Learning Specialization}
    University of Alberta on Coursera
}

% SHAD
\cvitem{
    \cvdurationstyle{2013 -- 2015}
}{
    \cvtitle{Diploma}
    Yandex School of Data Analysis
}

% uOttawa
\cvitem{
    \cvdurationstyle{2011 -- 2013}
}{
    \cvtitle{Master's degree, Systems Science}
    University of Ottawa
}

% uWaterloo
\cvitem{
    \cvdurationstyle{2006 -- 2010}
}{
    \cvtitle{Bachelor's degree, Mathematics / Business Administration}
    University of Waterloo
}

\end{document}
